\IfFileExists{stacks-project.cls}{%
\documentclass{stacks-project}
}{%
\documentclass{amsart}
}
\usepackage{amssymb}

% For dealing with references we use the comment environment
\usepackage{verbatim}
\newenvironment{reference}{\comment}{\endcomment}
%\newenvironment{reference}{}{}
\newenvironment{slogan}{\comment}{\endcomment}
\newenvironment{history}{\comment}{\endcomment}

% For commutative diagrams we use Xy-pic
\usepackage[all]{xy}

% We use 2cell for 2-commutative diagrams.
\xyoption{2cell}
\UseAllTwocells

% We use multicol for the list of chapters between chapters
\usepackage{multicol}

% This is generally recommended for better output
\usepackage{lmodern}
\usepackage[T1]{fontenc}
\usepackage{textalpha}

% For cross-file-references
\usepackage{xr-hyper}
\newcommand{\maybeexternaldocument}[2]{%
  \IfFileExists{#2.aux}{%
    \externaldocument[#1]{#2}%
  }{%
    \PackageWarning{stack}{Missing #2.aux; external references with prefix #1 unavailable}%
  }%
}
\let\externaldocumentorig\externaldocument
\renewcommand{\externaldocument}[2][]{%
  \IfFileExists{#2.aux}{%
    \externaldocumentorig[#1]{#2}%
  }{%
    \PackageWarning{stack}{Missing #2.aux; external references with prefix #1 unavailable}%
  }%
}

% Package for hypertext links:
\usepackage{hyperref}

% For any local file, say "hello.tex" you want to link to please
% use \externaldocument[hello-]{hello}
\externaldocument[introduction-]{introduction}
\externaldocument[conventions-]{conventions}
\externaldocument[sets-]{sets}
\externaldocument[categories-]{categories}
\externaldocument[topology-]{topology}
\externaldocument[sheaves-]{sheaves}
\externaldocument[sites-]{sites}
\externaldocument[stacks-]{stacks}
\externaldocument[fields-]{fields}
\externaldocument[algebra-]{algebra}
\externaldocument[brauer-]{brauer}
\externaldocument[homology-]{homology}
\externaldocument[derived-]{derived}
\externaldocument[simplicial-]{simplicial}
\externaldocument[more-algebra-]{more-algebra}
\externaldocument[smoothing-]{smoothing}
\externaldocument[modules-]{modules}
\externaldocument[sites-modules-]{sites-modules}
\externaldocument[injectives-]{injectives}
\externaldocument[cohomology-]{cohomology}
\externaldocument[sites-cohomology-]{sites-cohomology}
\externaldocument[dga-]{dga}
\externaldocument[dpa-]{dpa}
\externaldocument[sdga-]{sdga}
\externaldocument[hypercovering-]{hypercovering}
\externaldocument[schemes-]{schemes}
\externaldocument[constructions-]{constructions}
\externaldocument[properties-]{properties}
\externaldocument[morphisms-]{morphisms}
\externaldocument[coherent-]{coherent}
\externaldocument[divisors-]{divisors}
\externaldocument[limits-]{limits}
\externaldocument[varieties-]{varieties}
\externaldocument[topologies-]{topologies}
\externaldocument[descent-]{descent}
\externaldocument[perfect-]{perfect}
\externaldocument[more-morphisms-]{more-morphisms}
\externaldocument[flat-]{flat}
\externaldocument[groupoids-]{groupoids}
\externaldocument[more-groupoids-]{more-groupoids}
\externaldocument[etale-]{etale}
\externaldocument[chow-]{chow}
\externaldocument[intersection-]{intersection}
\externaldocument[pic-]{pic}
\externaldocument[weil-]{weil}
\externaldocument[adequate-]{adequate}
\externaldocument[dualizing-]{dualizing}
\externaldocument[duality-]{duality}
\externaldocument[discriminant-]{discriminant}
\externaldocument[derham-]{derham}
\externaldocument[local-cohomology-]{local-cohomology}
\externaldocument[algebraization-]{algebraization}
\externaldocument[curves-]{curves}
\externaldocument[resolve-]{resolve}
\externaldocument[models-]{models}
\externaldocument[functors-]{functors}
\externaldocument[equiv-]{equiv}
\externaldocument[pione-]{pione}
\externaldocument[etale-cohomology-]{etale-cohomology}
\externaldocument[proetale-]{proetale}
\externaldocument[relative-cycles-]{relative-cycles}
\externaldocument[more-etale-]{more-etale}
\externaldocument[trace-]{trace}
\externaldocument[crystalline-]{crystalline}
\externaldocument[spaces-]{spaces}
\externaldocument[spaces-properties-]{spaces-properties}
\externaldocument[spaces-morphisms-]{spaces-morphisms}
\externaldocument[decent-spaces-]{decent-spaces}
\externaldocument[spaces-cohomology-]{spaces-cohomology}
\externaldocument[spaces-limits-]{spaces-limits}
\externaldocument[spaces-divisors-]{spaces-divisors}
\externaldocument[spaces-over-fields-]{spaces-over-fields}
\externaldocument[spaces-topologies-]{spaces-topologies}
\externaldocument[spaces-descent-]{spaces-descent}
\externaldocument[spaces-perfect-]{spaces-perfect}
\externaldocument[spaces-more-morphisms-]{spaces-more-morphisms}
\externaldocument[spaces-flat-]{spaces-flat}
\externaldocument[spaces-groupoids-]{spaces-groupoids}
\externaldocument[spaces-more-groupoids-]{spaces-more-groupoids}
\externaldocument[bootstrap-]{bootstrap}
\externaldocument[spaces-pushouts-]{spaces-pushouts}
\externaldocument[spaces-chow-]{spaces-chow}
\externaldocument[groupoids-quotients-]{groupoids-quotients}
\externaldocument[spaces-more-cohomology-]{spaces-more-cohomology}
\externaldocument[spaces-simplicial-]{spaces-simplicial}
\externaldocument[spaces-duality-]{spaces-duality}
\externaldocument[formal-spaces-]{formal-spaces}
\externaldocument[restricted-]{restricted}
\externaldocument[spaces-resolve-]{spaces-resolve}
\externaldocument[formal-defos-]{formal-defos}
\externaldocument[defos-]{defos}
\externaldocument[cotangent-]{cotangent}
\externaldocument[examples-defos-]{examples-defos}
\externaldocument[algebraic-]{algebraic}
\externaldocument[examples-stacks-]{examples-stacks}
\externaldocument[stacks-sheaves-]{stacks-sheaves}
\externaldocument[criteria-]{criteria}
\externaldocument[artin-]{artin}
\externaldocument[quot-]{quot}
\externaldocument[stacks-properties-]{stacks-properties}
\externaldocument[stacks-morphisms-]{stacks-morphisms}
\externaldocument[stacks-limits-]{stacks-limits}
\externaldocument[stacks-cohomology-]{stacks-cohomology}
\externaldocument[stacks-perfect-]{stacks-perfect}
\externaldocument[stacks-introduction-]{stacks-introduction}
\externaldocument[stacks-more-morphisms-]{stacks-more-morphisms}
\externaldocument[stacks-geometry-]{stacks-geometry}
\externaldocument[moduli-]{moduli}
\externaldocument[moduli-curves-]{moduli-curves}
\externaldocument[examples-]{examples}
\externaldocument[exercises-]{exercises}
\externaldocument[guide-]{guide}
\externaldocument[desirables-]{desirables}
\externaldocument[coding-]{coding}
\externaldocument[obsolete-]{obsolete}
\externaldocument[fdl-]{fdl}
\externaldocument[index-]{index}

% Heap Project documents
\externaldocument[linear-algebra-]{linear-algebra}
\externaldocument[groups-]{groups}
\externaldocument[complex-analysis-]{complex-analysis}
\externaldocument[functional-analysis-]{functional-analysis}
\externaldocument[categories-]{categories}
\externaldocument[commutative-algebra-]{commutative-algebra}
\externaldocument[ringed-spaces-]{ringed-spaces}
\externaldocument[homological-algebra-]{homological-algebra}
\externaldocument[lie-algebras-]{lie-algebras}
\externaldocument[representation-theory-]{representation-theory}
\externaldocument[smooth-manifolds-]{smooth-manifolds}
\externaldocument[differential-topology-]{differential-topology}
\externaldocument[algebraic-topology-]{algebraic-topology}
\externaldocument[differential-geometry-]{differential-geometry}
\externaldocument[algebraic-geometry-]{algebraic-geometry}
\externaldocument[symplectic-geometry-]{symplectic-geometry}
\externaldocument[complex-geometry-]{complex-geometry}
\externaldocument[hodge-theory-]{hodge-theory}
\maybeexternaldocument{deformations-}{deformations}
\externaldocument[vertex-algebras-]{vertex-algebras}
\externaldocument[springer-]{springer}
\externaldocument[infty-categories-]{infty-categories}
\externaldocument[seminars-differential-geometry-]{%
seminars-differential-geometry}
\externaldocument[seminars-complex-geometry-]{%
seminars-complex-geometry}
\externaldocument[seminars-algebraic-geometry-]{%
seminars-algebraic-geometry}
\externaldocument[seminars-others-]{%
seminars-others}
%\externaldocument[geometric-prequantization-]{geometric-prequantization}
\externaldocument[surfaces-]{surfaces}
\externaldocument[birational-maps-commutative-algebra-]{birational-maps-commutative-algebra}
\externaldocument[cremona-transformations-]{cremona-transformations}
\externaldocument[derived-categories-of-sheaves-]{derived-categories-of-sheaves}
\externaldocument[geometric-invariant-theory-]{geometric-invariant-theory}
\externaldocument[geometric-stability-conditions-and-group-actions-]{geometric-stability-conditions-and-group-actions}
\externaldocument[moduli-spaces-of-sheaves-]{moduli-spaces-of-sheaves}
\externaldocument[bridgeland-stability-general-theory-]{bridgeland-stability-general-theory}
\externaldocument[hilbert-schemes-]{hilbert-schemes}
\externaldocument[noncommutative-abelian-surfaces-and-kummer-type-hyperkahler-varieties-]{noncommutative-abelian-surfaces-and-kummer-type-hyperkahler-varieties}
\externaldocument[mumford-tate-groups-in-hodge-theory-]{mumford-tate-groups-in-hodge-theory}
\externaldocument[hodge-birational-atoms-2026-]{hodge-birational-atoms-2026}
\externaldocument[xxii-egd-2026-]{%
xxii-egd-2026}
\externaldocument[pde-]{pde}
\externaldocument[probability-]{probability}

\externaldocument[linguistics-]{linguistics}
\externaldocument[genealogia-familiar-]{%
genealogia-familiar}

\externaldocument[%
16th-alga-meeting-2026-]{%
16th-alga-meeting-2026}

\externaldocument[%
iv-jornadas-lefschetz-2026-]{%
iv-jornadas-lefschetz-2026}

\externaldocument[reading-the-masters-]{%
reading-the-masters}

\externaldocument[real-analysis-]{%
real-analysis}

\externaldocument[optimal-transport-]{%
optimal-transport}

\externaldocument[history-]{history}

% Theorem environments.
%
\theoremstyle{plain}
\newtheorem{theorem}[subsection]{Theorem}
\newtheorem{proposition}[subsection]{Proposition}
\newtheorem{lemma}[subsection]{Lemma}

\theoremstyle{definition}
\newtheorem{definition}[subsection]{Definition}
\newtheorem{example}[subsection]{Example}
\newtheorem{exercise}[subsection]{Exercise}
\newtheorem{situation}[subsection]{Situation}

\theoremstyle{remark}
\newtheorem{remark}[subsection]{Remark}
\newtheorem{remarks}[subsection]{Remarks}

\numberwithin{equation}{subsection}

% Macros
%
\def\lim{\mathop{\mathrm{lim}}\nolimits}
\def\colim{\mathop{\mathrm{colim}}\nolimits}
\def\Spec{\mathop{\mathrm{Spec}}}
\def\Hom{\mathop{\mathrm{Hom}}\nolimits}
\def\Ext{\mathop{\mathrm{Ext}}\nolimits}
\def\SheafHom{\mathop{\mathcal{H}\!\mathit{om}}\nolimits}
\def\SheafExt{\mathop{\mathcal{E}\!\mathit{xt}}\nolimits}
\def\Sch{\mathit{Sch}}
\def\Mor{\mathop{\mathrm{Mor}}\nolimits}
\def\Ob{\mathop{\mathrm{Ob}}\nolimits}
\def\Sh{\mathop{\mathit{Sh}}\nolimits}
\def\NL{\mathop{N\!L}\nolimits}
\def\CH{\mathop{\mathrm{CH}}\nolimits}
\def\proetale{{pro\text{-}\acute{e}tale}}
\def\etale{{\acute{e}tale}}
\def\QCoh{\mathit{QCoh}}
\def\Ker{\mathop{\mathrm{Ker}}}
\def\Im{\mathop{\mathrm{Im}}}
\def\Coker{\mathop{\mathrm{Coker}}}
\def\Coim{\mathop{\mathrm{Coim}}}

% Boxtimes
%
\DeclareMathSymbol{\boxtimes}{\mathbin}{AMSa}{"02}

%
% Macros for moduli stacks/spaces
%
\def\QCohstack{\mathcal{QC}\!\mathit{oh}}
\def\Cohstack{\mathcal{C}\!\mathit{oh}}
\def\Spacesstack{\mathcal{S}\!\mathit{paces}}
\def\Quotfunctor{\mathrm{Quot}}
\def\Hilbfunctor{\mathrm{Hilb}}
\def\Curvesstack{\mathcal{C}\!\mathit{urves}}
\def\Polarizedstack{\mathcal{P}\!\mathit{olarized}}
\def\Complexesstack{\mathcal{C}\!\mathit{omplexes}}
% \Pic is the operator that assigns to X its picard group, usage \Pic(X)
% \Picardstack_{X/B} denotes the Picard stack of X over B
% \Picardfunctor_{X/B} denotes the Picard functor of X over B
\def\Pic{\mathop{\mathrm{Pic}}\nolimits}
\def\Picardstack{\mathcal{P}\!\mathit{ic}}
\def\Picardfunctor{\mathrm{Pic}}
\def\Deformationcategory{\mathcal{D}\!\mathit{ef}}


\begin{document}

\title{Afro-Brazilian art}
\maketitle

\phantomsection
\label{section-phantom}

\tableofcontents

\section{Method}
\label{section-afro-art-method}

This notebook began with the MUNCAB and
Casa do Benin visits of September 26,
2026.

The detailed order of photographed works,
labels, and questions is preserved in
\texttt{salvador-fieldnotes.tex}.

The main methodological distinction is
between a work, its museum label, a
curatorial interpretation, and a later
interpretation produced in conversation.
These layers should not be collapsed.

\section{MUNCAB as museum-territory}
\label{section-muncab-territory}

The Salvador itinerancy of the 36th Bienal
was titled
\emph{Nem todo viandante anda estradas---
Da humanidade como pr\'atica}.

The MUNCAB text called the institution a
``museum-territory''.  It asked how art
history makes presences, absences,
centralities, and silences, and insisted
that belonging does not require a
homogeneous identity
\cite{muncab2026bienal}.

The text explicitly placed its work in
dialogue with Concei\c{c}\~ao Evaristo,
L\'elia Gonzalez, Beatriz Nascimento,
Milton Santos, bell hooks, and
\'Edouard Glissant.

The point is not that the countries and
artists represented share one experience.
The curatorial problem is how historically
different territories were related by
colonialism, racialization, displacement,
memory, and artistic production.

\section{Hariel Revignet and water}
\label{section-hariel-water}

Hariel Revignet's exhibition was
\emph{Omi Tutu---\'Agua fresca}.  The
official MUNCAB text describes the artist
as Brazilian-Gabonese and makes water the
organizing language of territory,
continuity, return, and transformation
\cite{muncab2026omitutu}.

The most memorable curatorial image was a
river reaching the sea.  It can no longer
be distinguished, but it has not ceased to
exist.

This offers a model of diasporic continuity
without claiming that forms remain pure or
unchanged.  Memory persists through
mixture, displacement, and new channels.

The curatorial text also linked water with
Oxum, Iemanj\'a, and Nan\~a, and described
a feminine force that creates, sustains,
receives, and transforms.  This is a
curatorial reading of the works, not a
complete theology of those orix\'as.

\section{Ori Eja and Iya Eja}
\label{section-hariel-ori-iya}

A photographed work appeared beside the
labels \emph{Ori Ej\'a}, 2025, and
\emph{Iya Ej\'a}, 2025.

Both labels described acrylic on canvas
with cowrie-shell embroidery.  The single
photograph does not establish which label
belonged to the pictured canvas.  The
notebook must preserve that uncertainty.

The work showed a dark female figure, a
rounded abdomen, hands over the body, a
fish tail, and a surface densely worked
with cowries.

These visual observations justify
questions about fish, maternity, water,
and gestation.  They do not by themselves
justify identifying the figure as a
specific orix\'a.

\section{Ziquidas---Rio de Micangas}
\label{section-hariel-ziquidas}

The installation
\emph{Ziquid\'as---Rio de Mi\c{c}angas},
2026, suspended bead strings from woven
baskets above clay vessels, gourds, cowries,
seeds, bowls, and other containers.

The label named the materials in detail:
beads, gold-coloured brass wire,
a\c{c}a\'\i{} seeds, white clay vessels,
tiger cowries, a large clay vessel, gourds,
water jugs, a clay bowl, a basin, and
seed-fish.

The hanging strings could be seen as rain,
river, or falling water.  The connection
with waist beads and Gabonese family
memory is an interpretive and biographical
lead that requires a source before being
fixed as the work's meaning.

\section{Cowries and material meaning}
\label{section-art-cowries}

The works made cowries and beads actual
materials rather than painted signs.
Texture, weight, repetition, sound, labour,
and bodily association all enter the work.

The material can carry religious, economic,
ornamental, familial, and diasporic
histories.  No single symbolic gloss should
replace the work's material complexity.

\section{Diogum and the ofa}
\label{section-art-diogum}

A small iron sculpture beside the Diogum
biography had the form of a bow and central
arrow.  The museum text itself named the
ofa of Ox\'ossi among the artist's symbols.

Daniel's immediate recognition that the
object looked like Ox\'ossi rather than
Ogum was therefore correct.

The artist's name and work in iron remain
connected with Ogum.  The represented
attribute, however, can belong to Ox\'ossi.
This distinction prevents the artist's
identity from swallowing the formal content
of each work.

\section{Nadia Taquary and Iroko}
\label{section-art-nadia-iroko}

N\'adia Taquary's
\emph{Ir\^ok\^o: A \`arvore c\'osmica},
2025, used bead strings in colours
associated by the museum with Iroko.

The curatorial text joined the work to
ancestral time, the cycle of life, the
Brazilian gameleira, and the Iyami.

The work is important because religious
vocabulary enters a contemporary sculptural
language.  It should not be treated as a
neutral ethnographic diagram of Iroko.

The panel's translation of Iyami as
``feiticeiras'' also shows that museum
language can carry difficult acts of
translation.

\section{Sandro Aiye and Bara}
\label{section-art-sandro-bara}

Sandro Aiy\^e's
\emph{B\'ar\`a II}, 2025, was carved from
demolition wood and used beads.

The sculpture belonged to
\emph{Pad\^e On\~a---Encontrar Caminhos}.
MUNCAB's official account places Exu at
the centre through movement, mediation,
circulation, and the opening of paths
\cite{muncab2026padeona}.

The work should therefore be read through
its material, vertical presence, reused
wood, exhibition title, and relation to
Exu.  A trident-like form alone is not an
adequate interpretation.

\section{The sculpture garden}
\label{section-muncab-garden}

The MUNCAB sculpture garden occupies an
area formerly associated with a police
Delegacia de Jogos e Costumes.

The museum frames the conversion of that
space into a garden and exhibition area as
a reversal: a site connected with state
control of Black cultural practices becomes
a place for public art and encounter.

This institutional history belongs to the
meaning of the displayed works.  Site is
not merely background.

\section{Alberto Pitta}
\label{section-art-alberto-pitta}

Danillo Barata's
\emph{Alberto Pitta: F\'unf\'un D\'ud\'u}
collects and interprets Pitta's trajectory
through textile printing, painting, and
objects or installations
\cite{barata2025pitta}.

The key axis is carnival--terreiro--city.
Pitta's visual language shows how religious
and Afro-Bahian repertoires move into
fabric, clothing, procession, graphic
design, and urban public culture.

The book is therefore an art-historical
source.  It is not a general doctrinal
introduction to Candombl\'e.

\section{J. Cunha and Ile Aiye}
\label{section-art-jcunha}

\emph{J. Cunha e o Carnaval Negro}
addresses twenty-five years of J. Cunha's
visual work with Il\^e Aiy\^e
\cite{barata2024jcunha}.

This opens a study of costume, pattern,
colour, illustration, identity, and the
political visuality of the bloco afro.

Together, the Pitta and J. Cunha books show
that carnival is not outside serious art
history.  It is a major site of Black
visual production in Salvador.

\section{Zulu Araujo and first-person memory}
\label{section-art-zulu-araujo}

Zulu Ara\'ujo's
\emph{Apenas um cidad\~ao} is not primarily
an art monograph.  It belongs to memoir,
political history, and the history of Black
cultural institutions.

It remains relevant to art because Olodum,
public cultural policy, carnival, and Black
political organization shaped the field in
which Salvador's visual and musical culture
circulated.

The book should be used as first-person
memory, not silently converted into a
neutral comprehensive history.

\section{Tincoa and Os Tincoas}
\label{section-art-tincoa}

The exhibition \emph{Tinco\~a} used a
shared name to join the tinco\~a bird and
the Bahian vocal group Os Tinco\~as.

The curatorial movement was:

\begin{itemize}
\item bird call;
\item musical song;
\item ancestral voice;
\item memory in craft and material form.
\end{itemize}

This established music as part of the art
notebook.  Visual art, song, religion,
craft, and memory were being curated as
one cultural field rather than isolated
media.

\section{Casa do Benin as an artwork}
\label{section-art-casa-benin}

The Casa do Benin should not be understood
only as a room containing African objects.
The institution itself makes an argument
about the Bahia--Benin relation.

Pierre Verger's historical model of flux
and reflux provides the intellectual frame
\cite{verger2021fluxo}.  Lina Bo Bardi's
architectural intervention means that the
building also participates in this frame.

The museum thus joins collection,
architecture, history, and a reciprocal
institutional project between Salvador and
Ouidah.

\section{The vessel of freedom}
\label{section-art-vessel-freedom}

Emo de Medeiros's
\emph{Le vaisseau de la Libert\'e---O navio
da Liberdade}, 2025, showed an invented,
brightly patterned Atlantic vessel.

The image reverses or reworks the slave
ship as a symbol.  It should not be read as
a photograph of a historical ship.

The conversation proposed a connection to
Pierre Fatumbi Verger and to the artist's
\emph{Fatumbia} series.  That connection
should be retained as a lead requiring a
reliable catalogue source.

A claim that the image was produced with
generative AI was also made in conversation
but was not established by the photographed
label.  It should not enter the permanent
interpretation without verification.

\section{Mirongar and counter-visuality}
\label{section-art-mirongar}

The Casa do Benin exhibition text
\emph{Mirongar} described racism as a
history of discourse and image-making.

It proposed ``mirongar'' as a way to work
between what can be shown and what remains
guarded.  It linked images with routes of
escape, networks of affection, and a Black
counter-visuality.

This gives the visit an ethical question:
what is the difference between making
something visible and making it available
for consumption?

That question also applies to religious
images and rituals.  A museum's desire to
display does not cancel a community's
right to secrecy or restricted knowledge.

\section{Museum labels as objects}
\label{section-art-museum-labels}

The Akan mask was labelled as Central
African, while the museum's own map placed
the Akan region in Ghana and C\^ote d'Ivoire.

The contradiction should not be erased.
The label itself becomes evidence of how a
museum classifies and sometimes
misclassifies African objects.

Critical museum reading therefore includes
the frame, architecture, label, map,
curatorial text, and institutional history,
not only the isolated artwork.

\section{Research directions}
\label{section-afro-art-directions}

The next stages should include:

\begin{itemize}
\item a close reading of Hariel Revignet's
catalogue and artist statements;
\item the visual history of Il\^e Aiy\^e,
Olodum, and Cortejo Afro;
\item Alberto Pitta and textile as public
knowledge;
\item N\'adia Taquary's relation to Yoruba
religious and material traditions;
\item Black museum practice in Salvador;
\item Lina Bo Bardi's Casa do Benin;
\item the relation between secrecy,
religion, and museum display;
\item music, especially Os Tinco\~as, as a
companion archive to the visual material.
\end{itemize}

\bibliography{my,salvador}
\bibliographystyle{amsalpha}

\end{document}
